\documentclass[11pt,a4paper]{article}

\usepackage[margin=2.5cm]{geometry}
\usepackage[T1]{fontenc}
\usepackage{lmodern}
\usepackage{amsmath}
\usepackage{booktabs}
\usepackage{tabularx}
\usepackage{enumitem}
\usepackage{xcolor}
\usepackage{listings}
\usepackage{hyperref}

\hypersetup{colorlinks=true, linkcolor=blue!50!black, urlcolor=blue!50!black}
\lstset{
  basicstyle=\ttfamily\small,
  columns=fullflexible,
  breaklines=true,
  frame=single,
  framerule=0.3pt,
  rulecolor=\color{black!30},
  backgroundcolor=\color{black!3},
  keywordstyle=\color{blue!60!black}\bfseries,
  commentstyle=\color{black!55}\itshape,
  showstringspaces=false
}
\newcommand{\code}[1]{\texttt{#1}}

\title{CPU--GPU Co-Scheduling Module for CloudSim 7G\\[4pt]
\large Phase 1b Report: The GPU Layer (L1)}
\author{Chirag}
\date{21 November 2026}

\begin{document}
\maketitle

\begin{abstract}
Phase 1b builds the real GPU layer (L1) on CloudSim 7G and retires the stand-in written in Phase 1a.
The layer models GPU servers with gang-allocated whole GPUs, fractional CPU cores and memory, one
isolated guest per job part, training progress tracked directly in iterations, and server power from
CPU, GPUs and base load. Two timing problems in the earlier design were found and fixed: a wake-up
the CloudSim datacenter silently ignores, which made a sped-up job finish five times too late in a
test, and an ordering problem between a grow and a shrink on the same server in one round. The
randomized lock-step test from Check~1 now finishes 300 jobs with a worst error of 0.000000\,s
(Check~1: 0.037\,s). All 209 tests in the project pass.
\end{abstract}

\tableofcontents

%======================================================================
\section{Context and change of plan}
%======================================================================

The architecture assigned L1 to a teammate. In Phase~1a the contract between the scheduler and the
GPU layer (package \code{gpu.api}) was written together with a stub implementation, so the
scheduler could be developed first. On 21~November~2026 the author took over L1 as well. This
phase therefore implements L1 in full, in package \code{org.cloudbus.cloudsim.gpu} as named in the
architecture document, and removes the stub.

%======================================================================
\section{Design}
%======================================================================

\subsection{What the GPU layer does and does not decide}
The scheduler (L4) decides every job's servers, GPUs, cores, memory and speed. The GPU layer only
carries the decision out and reports state. In particular it never computes throughput: each
command carries the job's \emph{rate} in iterations per second, already including throughput and
the locality penalty. This is the ``rate-to-speed split'' agreed in Check~3.

\subsection{Classes}

\begin{center}
\begin{tabularx}{\linewidth}{@{}lX@{}}
\toprule
Class & Role \\
\midrule
\code{CloudSimGpuLayer} & Entry point (\code{GpuLayer}): builds one datacenter with one host per
  server, binds the executor to the broker, reports server status. \\
\code{GpuHost} & A server. Extends CloudSim's \code{PowerHost}. Holds a \code{GpuProvisioner},
  tracks cores and memory, enforces reservations, computes power. \\
\code{Gpu} & One whole, exclusive GPU: index, type, socket (NUMA hook, always 0), owner, busy fraction. \\
\code{GpuProvisioner} & Gang allocation: all requested GPUs at once, lowest-numbered first, or none. \\
\code{GpuVm} & The guest of one job part on one server (GPU passthrough isolation). \\
\code{TrainingCloudlet} & A part's work in iterations, as a \code{double}. \\
\code{GpuCloudletScheduler} & Advances work at the job's rate; ignores CloudSim's MIPS shares. \\
\code{TrainingJob} & A running job: its parts, which move in lock-step at one rate. \\
\code{GpuPlanExecutor} & Applies a round's commands (\code{Start}, \code{Resize}, \code{Stop}). \\
\code{LinearPowerModel} & $P(u) = P_\text{idle} + (P_\text{max} - P_\text{idle})\,u$, on at $u = 0$. \\
\code{PlannedGuestAllocationPolicy} & ``Do as told'' placement: a guest goes where the plan says or is refused. \\
\bottomrule
\end{tabularx}
\end{center}

\subsection{Progress in iterations, not MIPS}
CloudSim measures cloudlet work in integer millions of instructions and speed in MIPS shared by a
VM scheduler. The stub encoded speed as MIPS with a scale factor, which rounds work to integers and
ties speed to CPU shares that have no meaning for a GPU job. The architecture's fallback for this
case was ``track work directly in \code{GpuCloudletScheduler}'', and that is what the real layer does:

\begin{itemize}
  \item Each \code{TrainingCloudlet} holds its remaining work as a \code{double}.
  \item \code{GpuCloudletScheduler} keeps the job's rate $r$. On every update at time $t$ it subtracts
        $r\,(t - t_\text{prev})$ and predicts the finish at $t + \text{remaining}/r$.
  \item A rate change goes through \code{setRate(now, r')}, which first credits progress at the old
        rate. Contract rule~1 (credit before any change) is therefore enforced inside the scheduler.
  \item A cloudlet within $10^{-6}$\,s of done counts as done, absorbing floating-point noise at an
        exact finish event.
\end{itemize}

CloudSim's MIPS, RAM and bandwidth accounting is kept nominal (every guest asks for 1 unit) so it
never limits a job. The real resources are counted by \code{GpuHost}.

\subsection{Reservations happen when a command is applied}
CloudSim creates and destroys guests through events processed after the broker's turn. If GPUs,
cores and memory were taken at that point, a round that stops job~A and grows job~B on the same
server would see A's resources still held when B grows (A's destroy is still queued) and wrongly
report the server as full. The layer therefore reserves resources on \code{GpuHost} synchronously
when the executor applies a command. CloudSim's later creation event only carries the guest, and the
host accepts only guests that hold a reservation on it (contract rule~4).

\subsection{Applying a round in phases}
Within one \code{apply} call the executor works in three phases, so the order in which commands are
listed never matters:

\begin{enumerate}
  \item \textbf{Credit} all running jobs' progress at their old rates.
  \item \textbf{Release}: stops; moves (a resize to other servers or another GPU count is a stop
        followed by a start with the remaining work); and the old cores and memory of every
        in-place resize (the guest keeps its GPUs).
  \item \textbf{Reserve}: the new cores and memory of in-place resizes, then all starts.
\end{enumerate}

Phase~1a released and re-reserved each in-place resize one at a time. That fails when, on one
server with 24 cores, job~2 grows from 6 to 18 cores and is listed before job~1 shrinks from 18 to
6: the grow is checked while job~1 still holds 18 cores. The plan is feasible, but the order made it
look infeasible. Splitting \code{GpuHost}'s resize into \code{releaseCoresAndMemory} and
\code{reserveCoresAndMemory}, called in separate phases, fixes this. A test covers exactly this case.

\subsection{Waking the datacenter at the right time}
After a speed change, CloudSim's datacenter must recompute when the next cloudlet finishes (contract
rule~3). Phase~1a sent a ``wake now'' event. Reading \code{Datacenter.updateCloudletProcessing}
shows that the datacenter ignores an update requested within \code{minTimeBetweenEvents} (0.01\,s)
of its previous one:

\begin{lstlisting}[language=Java]
if (CloudSim.clock() < 0.111 ||
    CloudSim.clock() >= getLastProcessTime() + CloudSim.getMinTimeBetweenEvents()) { ... }
\end{lstlisting}

That is exactly the situation in event-driven scheduling, where a round runs at the moment a job
finishes and the datacenter has just processed that finish. The wake is dropped, and the only event
left is the one predicted from the old, slower rate. A new contract test reproduces it: job~1 ends at
100\,s and job~2 is sped up tenfold at the same instant. Job~2 should finish at 190\,s; under the
stub it finished at \textbf{1000\,s}.

\textbf{Fix.} After applying a round, the real executor schedules a datacenter update at each
started or resized job's new predicted finish time. That event is not within 0.01\,s of the
previous update, so it is processed, and the job finishes on time. Stale updates from earlier
predictions only cause harmless recomputations.

\paragraph{Remaining resolution limit.} The datacenter never schedules its next update sooner than
0.02\,s after the current one. Two jobs finishing less than 0.02\,s apart can have the second
reported up to 0.02\,s late. This is CloudSim's own time resolution. For jobs that run for hours it
is negligible, and it is stated here rather than worked around.

\subsection{Spread jobs}
A job over $k$ servers has $k$ guests. Each part's cloudlet holds the job's full remaining work and
runs at the job's rate, so all parts finish in the same datacenter update. The job ends, and the
listener is called, once the last part returns.

\subsection{Power}
Server power at any moment is
\[
  P = P_\text{CPU}(u_\text{CPU}) + \sum_{\text{GPUs}} P_\text{GPU}(u_\text{GPU}) + P_\text{base},
\]
with $u_\text{CPU}$ the allocated share of cores (allocated cores are assumed fully busy) and
$u_\text{GPU}$ each GPU's busy fraction. Commands carry a job's GPU busy fraction,
$\text{throughput}(c)/\text{throughput}(c^*)$, computed by the scheduler from the profiles. For
example, ResNet-18 at the fixed share runs at $1/2.4$ of full speed, so its GPUs are busy about
42\% of the time. Unallocated GPUs draw idle power: idle servers stay on.

\begin{itemize}
  \item The CPU model is either linear or one of CloudSim's SPECpower tables, chosen by class name.
  \item CloudSim's own \code{PowerModelLinear} returns 0\,W at zero utilisation, because it assumes
        an idle host is switched off. That contradicts the architecture, so \code{LinearPowerModel}
        is used instead. The SPECpower tables already report idle power at 0\%.
  \item \code{PowerSpec} has no default numbers. Every value must come from a stated source in the
        experiment configuration.
\end{itemize}

Power is piecewise constant between rounds and job ends. \code{GpuLayer.status(i)} returns each
server's used resources, utilisations and power, so L5 can integrate energy exactly by sampling at
those moments.

%======================================================================
\section{Contract changes}
%======================================================================
\begin{itemize}
  \item \code{JobCommand.Start} and \code{Resize} gained a GPU busy fraction in $[0, 1]$. The old
        four-argument forms remain and mean fully busy.
  \item \code{ServerShape} gained an optional \code{PowerSpec}; null means energy is not modelled.
  \item \code{GpuLayer} gained \code{status(i)}, returning a \code{ServerStatus}, and
        \code{serverCount()}.
  \item New records \code{PowerSpec} and \code{ServerStatus}.
\end{itemize}

%======================================================================
\section{Tests}
%======================================================================

The Phase~1a stub tests became an abstract \code{PlanExecutorContractTest}: the behaviour any
\code{GpuLayer} must have. Each implementation's test class extends it. The stub ran it too until
it was removed; it failed only the new wake-up test, which is the bug described above.

\subsection{Contract tests (9)}
\begin{center}
\begin{tabularx}{\linewidth}{@{}lX@{}}
\toprule
Test & Checks \\
\midrule
constant rate & 1000 iterations at 2/s finish at 500\,s; resources released. \\
in-place resize & Remaining 900 at 100\,s; finish at 325\,s after a 4$\times$ speed-up; cores updated. \\
stop and resume & Resume on another server with the remaining work; finish at 500\,s. \\
move & Resize to another server frees the old one. \\
same-round reuse & Stop and start in one round, start listed first. \\
over-commit & Nine GPUs on an 8-GPU server stops the run. \\
spread job & Two parts finish together; the job ends once. \\
speed-up at another job's end & Job~2 sped up at the instant job~1 ends finishes at 190\,s. \\
random resizes & 300 jobs, 1--4 parts, a new random rate every 300\,s; worst error 0.000000\,s. \\
\bottomrule
\end{tabularx}
\end{center}

\subsection{GPU-layer tests (7)}
\begin{center}
\begin{tabularx}{\linewidth}{@{}lX@{}}
\toprule
Test & Checks \\
\midrule
exact progress & Rate $1/3$: remaining work at 100\,s is $1000 - 100/3$ within $10^{-9}$. \\
grow and shrink & Grow listed before shrink on one 24-core server succeeds; both finish on time. \\
gang allocation & A 3-GPU and a 2-GPU job get GPUs 0--2 and 3--4; all freed at the end. \\
reservations only & A guest without a reservation is refused by the host. \\
power & Idle 600\,W; a 4-GPU ResNet-18-like job at 42\% busy and half the cores; then full speed,
  all cores. Each value matches the formula exactly, and CloudSim's \code{getPower()} agrees. \\
SPECpower & The IBM x3550 X5675 table gives 58.4\,W idle; an unknown name is rejected. \\
no power spec & Power reported as NaN. \\
\bottomrule
\end{tabularx}
\end{center}

\subsection{Whole project}
\begin{center}
\begin{tabular}{@{}lrr@{}}
\toprule
Module & Tests & Failures \\
\midrule
\code{cloudsim} (unchanged) & 167 & 0 \\
\code{cloudsim-examples} & 13 & 0 \\
\code{cloudsim-gpu} & 16 & 0 \\
\code{cloudsim-cosched} & 13 & 0 \\
\midrule
Total & 209 & 0 \\
\bottomrule
\end{tabular}
\end{center}
Command: \code{mvn clean install} from \code{cloudsim/}.

%======================================================================
\section{Assumptions and limits}
%======================================================================
\begin{itemize}
  \item Allocated CPU cores are fully busy; CPU utilisation is the allocated share.
  \item GPU busy fraction comes from the throughput profiles (computed by the scheduler).
  \item A linear GPU power model to rated maximum overestimates power; measured tables are preferred
        when available.
  \item Jobs on separate GPUs do not slow each other (no interference), as in Synergy.
  \item A move between servers has no cost in this version; restart cost belongs to the lifecycle
        policy (L4).
  \item NUMA and network links are not modelled; each GPU keeps a socket field as a hook.
  \item CloudSim's 0.02\,s minimum gap between datacenter updates bounds timing resolution.
\end{itemize}

%======================================================================
\section{Next steps}
%======================================================================
\begin{enumerate}
  \item File formats and loaders: job CSV, profile JSON, experiment YAML (\code{docs/formats.md}).
  \item L2 models: \code{ProfileRepository} and \code{ThroughputModel} from Synergy's profiles
        (\code{simulator/models/*.py}, now available), \code{LocalityModel}; fidelity test
        (ResNet-18 at 9 cores per GPU $= 2.4\times$).
  \item L3 workload: \code{TraceWorkload} and the labellers.
  \item Then Phase~2: the broker and the first end-to-end run.
\end{enumerate}

\end{document}
